\documentclass[10pt,a4paper]{amsart}
\usepackage[margin=19mm]{geometry}
\usepackage{amsmath,amssymb,mathtools}
\usepackage[scr=rsfs,cal=euler]{mathalfa}
\usepackage{xfrac}
\usepackage[unicode,hidelinks,bookmarks=false]{hyperref}
\newcommand{\ie}{i.e.\ }
\newcommand{\cf}{cf.\ }
\newcommand{\resp}{respectively\ }
\newcommand{\Subsection}{\subsection}
\providecommand{\nobreakdash}{\nobreak\hbox{-}}
\setlength{\emergencystretch}{2em}
\multlinegap0pt
\usepackage{amssymb}
\usepackage{float}
\newtheorem{theorem}{Theorem}
\title{Runs in Paperfolding Sequences}
\author{Jeffrey Shallit}
\date{Source: arXiv:2412.17930v2 (CC BY 4.0)}
\begin{document}
\maketitle

\noindent\emph{Source and licence.} Runs in Paperfolding Sequences. Version arXiv:2412.17930v2 (CC BY 4.0), distributed by arXiv under the Creative Commons Attribution 4.0 International licence, http://creativecommons.org/licenses/by/4.0/, as recorded in the arXiv licence metadata for this exact version and shown on the abstract page https://arxiv.org/abs/2412.17930v2. Full licence notice: \texttt{licenses/arxiv-2412.17930-CC-BY-4.0.txt}.\par\smallskip

\noindent\textit{Editorial scope.} Counted unit: the article's theorem theorem (source0.tex lines 710--720). The article's complete original proof of it is delivered with it (source0.tex lines 722--758, one \texttt{proof} environment of 37 lines). No mathematics is written, repaired or re-derived: the statement and the proof are the article's own lines, in the article's own notation, and no retained line is substituted or re-flowed.  No further source-local context is needed: every cross-reference of every retained line resolves inside the two delivered spans.  One declared adaptation: exactly 1 ancillary strings inside the retained spans are omitted (image-inclusion commands and, where the source repeats a label, the repeated label definitions) and no other line differs from the source; the figure environments, with their placement, captions and labels, are kept, and the package records the omitted strings in fidelity receipts bound to the affected bodies.  No retained line contains a citation, so the wrapper carries no bibliography and no bibliography entry.  The retained lines contain the article's own diagram code, which the wrapper typesets with the standard tikz packages; no image file is delivered and the package declares no asset dependency.  Editorial formatting only, in addition: an AMS \texttt{amsart} preamble that restates the article's own declarations and macros used by the retained lines, namely the environments \texttt{theorem} on separate counters, together with the standard packages those lines need.  The retained environments print in this document as Theorem 1 in order of appearance, whereas the article prints the same items with the numbering of its own full document.  The complete native source of the article as distributed in the arXiv e-print for this exact version is bundled unchanged as \texttt{sources/170632/source0.tex}.
\begin{theorem}
Arrange the set of positive integers
not in $H := \{ h(n)+1 \, : \, n \geq 0 \}$ in 
increasing order, and let $t(n)$ be the $n$'th such integer, for $n \geq 1$.
Then
\begin{itemize}
\item[(a)] $g(h(i)+1) = 2$ for $i \geq 0$;
\item[(b)] $g(t(2i)) = 3$ for $i \geq 1$;
\item[(c)] $g(t(2i-1)) = 1$ for $i \geq 1$.
\end{itemize}
\end{theorem}

\begin{proof}
The first step is to create an automaton {\tt tt} computing $t(n)$.  Once again,
we guess the automaton from data and then verify its correctness.
It is depicted in Figure~\ref{fig9}.
\begin{figure}[htb]
\begin{center}

\end{center}
\caption{The automaton {\tt tt} computing $t(n)$.}
\label{fig9}
\end{figure}

In order to verify its correctness,
we need to verify that {\tt tt} indeed computes a increasing function $t(n)$
and further that the set $\{ t(n) \, : \, n \geq 1 \} =
\{ 1,2, \ldots, \} \setminus H$. We can do this as follows:
\begin{verbatim}
eval tt1 "?lsd_2 An (n>=1) => Ex $tt(n,x)":
# takes a value for all n
eval tt2 "?lsd_2 ~En,x,y n>=1 & x!=y & $tt(n,x) & $tt(n,y)":
# does not take two different values for the same n
eval tt3 "?lsd_2 An,y,z (n>=1 & $tt(n,y) & $tt(n+1,z)) => y<z":
# is an increasing function
eval tt4 "?lsd_2 Ax (x>=1) => 
   ((En n>=1 & $tt(n,x)) <=> (~Em,y $ep_reg(m,y) & x=y+1))":
# takes all values not in H
\end{verbatim}

Now we can verify parts (a)-(c) as follows:
\begin{verbatim}
eval parta "?lsd_2 Ai,x (i>=1 & $ep_reg(i,x)) => RLR[x+1]=@2":
eval partb "?lsd_2 Ai,x (i>=1 & $tt(2*i,x)) => RLR[x]=@3":
eval partc "?lsd_2 Ai,x (i>=1 & $tt(2*i-1,x)) => RLR[x]=@1":
\end{verbatim}
And {\tt Walnut} returns {\tt TRUE} for all of these.  This completes the
proof.
\end{proof}

\end{document}
